\section{Estrutura do documento}
\label{sec:estrutura}

O arquivo \texttt{main.tex} coordena o projeto. A classe é carregada com
\verb|\documentclass{relatorioic}|, os dados da capa são alterados com
\verb|\renewcommand| e cada parte é inserida com \verb|\import|.

\begin{lstlisting}[language={[LaTeX]TeX},caption={Estrutura mínima do documento},label={cod:estrutura}]
\documentclass{relatorioic}
\renewcommand{\titulo}{Titulo do relatorio}
\renewcommand{\nome}{Nome do autor}
\renewcommand{\email}{autor@unesp.br}
\renewcommand{\orientador}{Nome do orientador}
\renewcommand{\emailorientador}{orientador@unesp.br}
\renewcommand{\cidade}{Presidente Prudente}
\renewcommand{\cidadeestado}{\cidade\;-- SP}
\addbibresource{referencias.bib}

\begin{document}
\import{sections/textual/}{6-introducao.tex}
\end{document}
\end{lstlisting}

\subsection{Seções e parágrafos}

Use os comandos abaixo para criar a hierarquia que aparecerá automaticamente no
sumário. Uma linha vazia no arquivo-fonte inicia um novo parágrafo.

\begin{lstlisting}[language={[LaTeX]TeX},caption={Hierarquia de títulos}]
\section{Titulo da secao}
\subsection{Titulo da subsecao}
\subsubsection{Titulo da subsubsecao}
\end{lstlisting}

\subsection{Sumário e listas automáticas}

As listas pré-textuais são geradas a partir dos títulos, legendas e siglas
cadastrados ao longo do documento. É preciso compilar mais de uma vez para que
todos os números de página sejam atualizados.

\begin{lstlisting}[language={[LaTeX]TeX},caption={Comandos para índices e listas}]
\tableofcontents       % Sumario
\listoffigures         % Lista de figuras
\listoftables          % Lista de tabelas
\lstlistoflistings     % Lista de codigos
\printnomenclature     % Lista de siglas
\end{lstlisting}

\subsection{Caracteres reservados}

Alguns caracteres têm função especial e precisam de uma barra invertida para
serem impressos. Por exemplo: \verb|\%| produz \%, \verb|\&| produz \&,
\verb|\_| produz \_ e \verb|\#| produz \#.
